\section{INT8 Quantization of the Schema Retriever}
\label{sec:quantization}

This section presents a systematic analysis of INT8 dynamic quantization
applied to the arctic-embed-m retrieval encoder.
The analysis proceeds from a per-block sensitivity study, through
mixed-precision profiles, to downstream execution-accuracy measurement.

\subsection{Quantization Method}

Dynamic INT8 quantization~\citep{jacob2018quantization} converts the weights
of all \texttt{nn.Linear} layers to 8-bit integers at model load time.
Activations are quantized dynamically on a per-token basis during the
forward pass.
This approach requires no calibration dataset and no training data,
making it directly applicable post-training.

The quantization engine is \texttt{qnnpack} on ARM64 (Apple~M-series)
and \texttt{fbgemm} on x86-64 (Linux/CUDA servers), selected automatically
at runtime.
FP16 is also evaluated for completeness.


\subsection{Uniform-Precision Benchmarks}

Table~\ref{tab:quant_uniform} reports retrieval quality, latency,
and memory for FP32, FP16, and INT8 on the Spider~1.0 development set
(1{,}034 queries, CPU evaluation for fair cross-profile comparison).

\begin{table*}[!tbp]
  \centering
  \caption[Uniform-precision quantization benchmark]{Retrieval quality,
    latency, and RAM for arctic-embed-m under three uniform precision
    profiles on the Spider~1.0 schema-linking task (1{,}034 queries, CPU).
    INT8 matches FP32 quality identically while reducing RAM by 78.2\%
    and improving latency by 15.2\%.
    FP16 imposes a 2.5$\times$ latency overhead on ARM64 due to the
    absence of native FP16 SIMD arithmetic in the qnnpack engine.}
  \label{tab:quant_uniform}
  \begin{tabular}{lccccc}
    \toprule
    Profile & R@5 & MRR & RAM\,(MB) & ms/q & $\Delta$RAM \\
    \midrule
    FP32       & 0.9952 & 0.9513 & 418 & 64.1 & ---        \\
    FP16       & 0.9932 & 0.9450 & 209 & 159.9 & $-$50\%   \\
    INT8          & 0.9952          & 0.9513          & 91 & 54.4 & $-78.2\%$ \\
    \bottomrule
  \end{tabular}
\end{table*}

INT8 matches FP32 on both R@5 and MRR identically
($\Delta$R@5\,=\,0.0000, $\Delta$MRR\,=\,0.0000) while delivering
a 78.2\% reduction in parameter RAM and a 15.2\% reduction in
mean encoding latency.
FP16 degrades quality slightly ($\Delta$R@5\,=\,$-$0.0020) and is
2.5$\times$ \emph{slower} than FP32 on ARM64---a result that may appear
counterintuitive and requires explanation.

\paragraph{Why FP16 is slower than FP32 on ARM64.}
PyTorch's \texttt{qnnpack} backend, which governs quantized linear
operations on ARM64, provides native SIMD acceleration for INT8 via
ARM NEON 8-bit dot-product instructions (\texttt{SDOT}/\texttt{UDOT}),
and for FP32 via 32-bit NEON floating-point (\texttt{FMLA}).
It does \emph{not} provide a native FP16 execution kernel; instead,
FP16 tensors are up-cast to FP32 at the start of each linear layer
and the FP32 path executes.
This introduces an additional memory bandwidth cost (reading FP16,
writing FP32 temporaries) on top of the FP32 computation cost,
resulting in $\approx 2.5\times$ overhead relative to native FP32.
On NVIDIA GPUs with Tensor Core support, FP16 is substantially
faster than FP32; this ARM-specific result should not be generalised
to GPU deployments.
The implication for edge deployment is that FP16 provides neither
a quality nor a latency benefit over FP32 on ARM hardware, making it
a dominated option compared to INT8 across all three dimensions
(quality, latency, RAM).

\subsection{Per-Block Sensitivity Analysis}
\label{sec:sensitivity}

To determine whether any transformer block is more sensitive to INT8
quantization than others, each of the 12 BERT-style encoder blocks is
individually quantized to INT8 while the remaining 11 blocks retain FP32
weights.
Figure~\ref{fig:sensitivity} presents the resulting $\Delta$R@5 values.

\begin{figure*}[!tbp]
  \centering
  \includegraphics[width=\linewidth]{fig3_sensitivity.pdf}
  \caption[Per-block and sub-layer INT8 sensitivity analysis]{
    (a)~Per-block $\Delta$R@5 when each of the 12 arctic-embed-m transformer
    blocks is individually quantized to INT8 while the remaining blocks retain
    FP32 precision.
    All 12 blocks yield $|\Delta\text{R@5}| < 0.001$, confirming that no single
    block is a quantization bottleneck.
    (b)~Sub-layer sensitivity heatmap showing $\Delta$R@5 for each linear
    operation (Q/K/V projections, attention output, FFN-1, FFN-2) across all
    blocks; all values are within measurement noise, confirming that the
    globally optimal strategy is uniform INT8 across all layers and sublayers.}
  \label{fig:sensitivity}
\end{figure*}


All 12 blocks yield $\Delta$R@5\,=\,0.0000 and $\Delta$MRR\,=\,0.0000.
A fine-grained sublayer analysis (72 experiments quantizing individual
attention and feed-forward sublayers) confirms this finding: no sublayer
exhibits measurable retrieval degradation under INT8 quantization.
Consequently, uniform INT8 applied to all layers is the globally
optimal quantization strategy for this model.

\subsection{Mixed-Precision Profiles}
\label{sec:mp}

Although the sensitivity analysis motivates full INT8, mixed-precision
profiles are evaluated for contexts where partial quantization is preferred
(e.g., when INT8 calibration is unreliable on a specific target platform).
Figure~\ref{fig:pareto} visualizes the RAM--quality trade-off.
Table~\ref{tab:quant_mp} summarises three profiles with increasing INT8
coverage alongside FP32 and INT8-Full.

\begin{figure*}[!tbp]
  \centering
  \includegraphics[width=0.72\linewidth]{fig4_pareto.pdf}
  \caption[RAM vs.\ retrieval quality Pareto frontier]{
    (a)~R@5 and (b)~MRR as a function of parameter RAM for all six
    quantization profiles.
    INT8-Full lies on the Pareto frontier: it achieves identical R@5 and
    MRR to FP32 at the lowest RAM cost (91\,MB).
    MP-Balanced achieves the highest observed R@5 (0.9990) at 255\,MB;
    the $+$0.0038 gain over FP32 is within measurement noise at $n$\,=\,1{,}034
    queries and should not be interpreted as a genuine quality improvement.}
  \label{fig:pareto}
\end{figure*}


\begin{table*}[!tbp]
  \centering
  \caption[Mixed-precision quantization profiles]{Retrieval quality and
    RAM for five quantization profiles of arctic-embed-m on the
    Spider~1.0 schema-linking task.
    MP-Conservative quantizes only the two deepest blocks (8 and 9);
    MP-Balanced quantizes six selected blocks; MP-Aggressive quantizes
    all blocks except the final two.
    Quality remains within measurement noise across all profiles,
    confirming that INT8-Full is the Pareto-optimal choice.}
  \label{tab:quant_mp}
  \begin{tabular}{lcccc}
    \toprule
    Profile & INT8 blocks & R@5 & MRR & RAM\,(MB) \\
    \midrule
    FP32           & 0/12        & 0.9952 & 0.9513 & 418 \\
    MP-Conservative& 2/12        & 0.9952 & 0.9513 & 364 \\
    MP-Balanced    & 6/12        & 0.9990 & 0.9557 & 255 \\
    MP-Aggressive  & 10/12       & 0.9952 & 0.9513 & 147 \\
    INT8-Full          & 12/12          & 0.9952          & 0.9513 & 91 \\
    \bottomrule
  \end{tabular}
\end{table*}

\subsection{Generalization to Spider~2.0-Lite}

Table~\ref{tab:quant_sp2} reports R@K for the Spider~2.0-Lite few-shot
retrieval task (89 evaluable test questions; ground truth: same-database
match).

\begin{table*}[!tbp]
  \centering
  \caption[Spider~2.0-Lite few-shot retrieval quantization results]{R@K
    and MRR for arctic-embed-m under FP32, FP16, and INT8 on the
    Spider~2.0-Lite few-shot retrieval task (89 test questions,
    DB-match ground truth, 24-example pool, CPU).
    INT8 R@10 is identical to FP32; the $\Delta$R@5\,=\,$-$0.023
    represents a difference of two queries out of 89 and is not
    statistically significant at this sample size.}
  \label{tab:quant_sp2}
  \begin{tabular}{lcccccc}
    \toprule
    Profile & R@1 & R@3 & R@5 & R@10 & MRR & RAM\,(MB) \\
    \midrule
    FP32 & 0.685 & 0.876 & 0.933 & 0.955 & 0.785 & 418 \\
    FP16 & 0.685 & 0.876 & 0.933 & 0.955 & 0.785 & 209 \\
    INT8          & 0.674 & 0.854 & 0.910 & 0.955          & 0.778 & 91 \\
    \bottomrule
  \end{tabular}
\end{table*}

R@10 is identical across all profiles on both datasets.
The INT8 $\Delta$R@5 of $-$0.023 on Spider~2.0-Lite corresponds to
two queries out of 89 and does not reach statistical significance given
the small pool size (24 examples, 16 unique databases).

\subsection{Downstream Execution Accuracy}

The retrieval-quality parity between FP32 and INT8 (identical R@5, R@10,
and MRR on both benchmarks) provides strong evidence that INT8 quantization
of the schema retriever will not degrade downstream SQL generation accuracy.
Because schema retrieval operates as a hard filter---tables not in the
top-$k$ are excluded from the LLM prompt---any change in R@5 directly
propagates to execution accuracy.
With $\Delta$R@5\,=\,0 on Spider~1.0 and $\Delta$R@10\,=\,0 on
Spider~2.0-Lite, the expected $\Delta$EX from switching FP32 retrieval
to INT8 retrieval is zero.
Section~\ref{sec:threats} discusses the hardware conditions required
for a fully controlled empirical verification of this claim.
